\documentclass[11pt]{article}
\usepackage[a4paper,margin=27mm]{geometry}
\usepackage{amsmath,amssymb,amsthm,booktabs,array}
\usepackage[T1]{fontenc}
\usepackage{lmodern}
\usepackage[colorlinks=true,linkcolor=blue,citecolor=blue,urlcolor=blue]{hyperref}
\newcommand{\U}{\mathrm U}
\newcommand{\tr}{\operatorname{tr}}
\newcommand{\diag}{\operatorname{diag}}
\newcommand{\op}{\mathrm{op}}
\newcommand{\HS}{\mathrm{HS}}
\newtheorem{theorem}{Theorem}
\newtheorem{corollary}[theorem]{Corollary}
\title{Trace-coupled counterexamples to Larotonda's\texorpdfstring{\\}{} principal-log distance conjecture:\texorpdfstring{\\}{} beyond the Antezana--Larotonda--Varela theorem}
\author{Haobo Ma\\\small ChronoAI Pte Ltd, Singapore
\and Wenlin Zhang\\\small National University of Singapore\\\small ChronoAI Pte Ltd, Singapore}
\date{16 November 2026}
\begin{document}
\maketitle
\begin{abstract}
Beyond the unitarily invariant theorem of Antezana, Larotonda and Varela
(2011/2014; \cite[Theorem~4.3]{ALV}), Larotonda's conjecture
after Proposition 4.3 in Section 4.1, p.~7 of
\emph{Short curves in homogeneous spaces of the unitary group for unitarily
invariant norms} (arXiv:2610.31535v1) fails for a normalized, symmetric
Ad-invariant quadratic norm on $\mathfrak u(2)$: shifting one eigenangle
by $2\pi$ preserves the endpoint and gives a shorter exponential curve.
For the full positive-definite family obtained by adding $c$ times
the squared trace to the Hilbert--Schmidt norm squared, we show in every
dimension $n\geq2$ that principal-log minimality holds for all
principal angles if and only if the trace coefficient is zero.
Scalar examples give a strict failure whenever their eigenangle exceeds
$\pi(1+c)/(1+nc)$, for $c>0$.
Thus the proposed $\pi/2$ input already fails in dimension three,
and no positive dimension-independent spectral radius suffices.
\end{abstract}

\section{The question and the quadratic family}
Write $\U(n)=\{u\in M_n(\mathbb C):u^*u=I\}$ and
$\mathfrak u(n)=\{z\in M_n(\mathbb C):z^*=-z\}$.
For a symmetric norm $F$ on this real vector space, assume
\[
 F(uzu^*)=F(z)\quad(u\in\U(n)),\qquad \|z\|_{\op}\leq F(z).
\]
Here $\|\cdot\|_{\op}$ is the Euclidean operator norm.
For a piecewise smooth curve $\gamma:[0,1]\to\U(n)$ put
\[
 L_F(\gamma)=\int_0^1F(\gamma(t)^*\dot\gamma(t))\,dt,
 \qquad d_F(u,v)=\inf_{\gamma(0)=u,\,\gamma(1)=v}L_F(\gamma).
\]
The conjecture in \cite[Section 4.1, p.~7, after Proposition 4.3]{Larotonda}
is refuted below, even within this class of symmetric norms. It asserts
\begin{equation}\label{eq:conjecture}
 d_F(I,e^z)=F(z)\qquad (z\in\mathfrak u(n),\ \|z\|_{\op}\leq\pi).
\end{equation}
Its stated class also allows nonsymmetric Finsler norms, so a symmetric
counterexample suffices.

The distinction from singular-value invariance is essential.
Antezana, Larotonda and Varela prove \eqref{eq:conjecture} for unitarily
invariant norms \cite[preprint Theorem 4.3]{ALV}.
Russell and Stepney discuss minimization over logarithm branches in a
quantum-control setting and the sufficient role of eigenvalue-modulus
monotonicity \cite[Section 4, after equation (11)]{RS}.
Here the obstruction is the trace, which depends on signed eigenangles,
not just singular values. Local minimizing-segment results for general
Ad-invariant norms, such as \cite[Theorem 4.11]{Metric} and
\cite[Corollary 2.17]{LR}, concern a ball in
the chosen Finsler norm and do not assert the spectral-radius bound
in \eqref{eq:conjecture}.

The eigenangle-convexity results of Larotonda and Miglioli
\cite[Theorems 4.2 and 5.9]{LM} supply a different ingredient:
convexity along segments whose principal logarithms remain in the
spectral ball of radius $\pi/2$.

Let $n\geq1$ and $c>-1/n$. For Hermitian $H$ define
\begin{equation}\label{eq:norm}
 Q_c(H)=\tr(H^2)+c(\tr H)^2,\qquad
 s_{n,c}=\frac1{\sqrt{\min\{1,1+nc\}}},\qquad
 F_c(iH)=s_{n,c}\sqrt{Q_c(H)}.
\end{equation}
Every skew-Hermitian matrix has a unique expression $iH$ with $H$ Hermitian.
The Hilbert--Schmidt norm is $\|H\|_{\HS}=\sqrt{\tr(H^2)}$.
Writing $H_0=H-(\tr H)I/n$, one obtains
\begin{equation}\label{eq:split}
 Q_c(H)=\tr(H_0^2)+(1/n+c)(\tr H)^2.
\end{equation}
This is a positive-definite real quadratic form precisely when $c>-1/n$.
It defines a norm by its associated inner product, and conjugation
preserves both terms in \eqref{eq:norm}.
For $c\geq0$, $Q_c(H)\geq\|H\|_{\HS}^2$.
For $c<0$, Cauchy--Schwarz gives $(\tr H)^2\leq n\tr(H^2)$, whence
$Q_c(H)\geq(1+nc)\|H\|_{\HS}^2$.
Thus $F_c(iH)\geq\|H\|_{\HS}\geq\|H\|_{\op}$ in both cases.
This verifies all the stated norm hypotheses.

\section{A shorter logarithm in every nonzero case}
\begin{theorem}\label{thm:family}
Let $n\geq2$ and $-1/n<c\ne0$. Define
\begin{gather*}
 \varepsilon=\frac{|c|}{4(1+c)},\quad
 a=\pi(1-\varepsilon),\quad b=\frac{\pi}{2}\operatorname{sgn}(c),\\
 H=\diag(a,b,0,\ldots,0),\quad
 H'=\diag(a-2\pi,b,0,\ldots,0).
\end{gather*}
Then $\|iH\|_{\op}<\pi$, $e^{iH'}=e^{iH}$, and
\begin{equation}\label{eq:strict}
 d_{F_c}(I,e^{iH})\leq F_c(iH')<F_c(iH).
\end{equation}
Consequently, within the family \eqref{eq:norm},
\eqref{eq:conjecture} holds for every admissible $z$ if and only if $c=0$.
\end{theorem}
\begin{proof}
The denominator $1+c$ is positive. For $c>0$,
$0<\varepsilon<1/4$. For $c<0$, put $d=-c$; then
$0<d<1/n\leq1/2$, so
$\varepsilon=d/(4(1-d))<1/4$ again.
Hence $3\pi/4<a<\pi$ and $|b|=\pi/2$, proving the operator-norm bound.
The two diagonal exponentials agree because their first eigenangles
differ by $2\pi$.

For arbitrary real angles $\theta_1,\ldots,\theta_n$, set
$T=\sum_j\theta_j$ and $D=\diag(\theta_1,\ldots,\theta_n)$.
Replacing $\theta_1$ by $\theta_1-2\pi$ gives a diagonal matrix $D'$ with
\begin{align}\label{eq:gap}
 Q_c(D')-Q_c(D)
 &=(-4\pi\theta_1+4\pi^2)
     +c(-4\pi T+4\pi^2)\notag\\
 &=4\pi\big((\pi-\theta_1)+c(\pi-T)\big).
\end{align}
For $D=H$, $T=a+b$, so the expression in parentheses is
\[
 (1+c)\pi\varepsilon-cb
 =\frac{\pi|c|}{4}-\frac{\pi|c|}{2}
 =-\frac{\pi|c|}{4}.
\]
It follows that $Q_c(H')-Q_c(H)=-\pi^2|c|<0$.
Positive scaling and taking square roots preserve this strict inequality.

The smooth unitary curve $\gamma(t)=e^{itH'}$ has
$\gamma(0)=I$, $\gamma(1)=e^{iH}$ and
$\gamma(t)^*\dot\gamma(t)=iH'$.
Its constant speed gives $L_{F_c}(\gamma)=F_c(iH')$, proving
\eqref{eq:strict}. Finally $c=0$ is exactly the Hilbert--Schmidt norm,
for which \eqref{eq:conjecture} is the cited theorem of
Antezana, Larotonda and Varela. This proves the asserted equivalence.
\end{proof}

For the concrete two-dimensional choice $c=4$ the scale in
\eqref{eq:norm} is $1$, and the theorem reads
\begin{equation}\label{eq:main-example}
 H=\pi\diag(4/5,1/2),\quad H'=\pi\diag(-6/5,1/2),\quad
 \begin{cases}
 F_4(iH)^2=153\pi^2/20,\\
 F_4(iH')^2=73\pi^2/20.
 \end{cases}
\end{equation}
Indeed the two Hilbert--Schmidt squared terms are $89\pi^2/100$
and $169\pi^2/100$, while the squared-trace contributions are
$676\pi^2/100$ and $196\pi^2/100$.
Thus the increased Hilbert--Schmidt term is more than offset by the
decrease in the trace term. For $c\ne0$, take two diagonal Hermitian
matrices whose first two entries are $(1,1)$ and $(1,-1)$, respectively,
and whose other entries are zero. They have the same singular values but
distinct $F_c$-norms after multiplication by $i$.
Thus $F_c$ is not unitarily invariant in the stronger sense
required in \cite{ALV}.

\section{Scalar examples and the smaller-radius input}
Section 4.1 of \cite{Larotonda} conditionally proposes using
principal-log minimality for $\|z\|_{\op}<\pi/2$ to extend its quotient
argument. This is not a separately asserted theorem.
The same winding identity shows that the proposed input fails.
\begin{corollary}\label{cor:scalar}
Let $n\geq2$, $c>0$, and let $P=\diag(1,0,\ldots,0)$.
For
\[
 \frac{\pi(1+c)}{1+nc}<\theta<\pi,\qquad
 H_\theta=\theta I,\qquad H_\theta'=H_\theta-2\pi P,
\]
one has $e^{iH_\theta'}=e^{iH_\theta}$ and
$d_{F_c}(I,e^{iH_\theta})<F_c(iH_\theta)$.
For every fixed $n\geq2$ and every $r>\pi/n$, there are $c>0$
and such a counterexample with $\|iH_\theta\|_{\op}<r$.
\end{corollary}
\begin{proof}
Formula \eqref{eq:gap}, with $T=n\theta$, gives
\[
 Q_c(H_\theta')-Q_c(H_\theta)
 =4\pi\big(\pi(1+c)-\theta(1+nc)\big)<0.
\]
The exponentials agree by periodicity; the shorter exponential curve
gives the distance inequality as in Theorem~\ref{thm:family}.
For the last assertion choose $\pi/n<\theta<\min\{r,\pi\}$ and then
\[
 c>\frac{\pi-\theta}{n\theta-\pi}>0.
\]
This is exactly the displayed strict inequality. At
$\theta=\pi(1+c)/(1+nc)$ the two displayed logarithms have equal norm;
no distance formula is inferred from that equality.
\end{proof}

\begin{corollary}\label{cor:half}
Principal-log minimality for the full normalized Ad-invariant class
fails even when $\|z\|_{\op}<\pi/2$. In dimension three take $c=4$ and
\[
 H=\pi\diag(2/5,2/5,2/5),\qquad H'=\pi\diag(-8/5,2/5,2/5).
\]
Then
\[
 \|iH\|_{\op}=2\pi/5<\pi/2,\qquad
 F_4(iH)^2=156\pi^2/25,\quad
 F_4(iH')^2=136\pi^2/25,
\]
and $d_{F_4}(I,e^{iH})<F_4(iH)$.
Moreover, no positive spectral radius independent of dimension makes
principal-log minimality valid for all the norms $F_4$.
\end{corollary}
\begin{proof}
The displayed values follow by summing the squares of the diagonal
entries and four times the square of their sum. Endpoint equality and
the shorter curve are as in Theorem~\ref{thm:family}.
For the last assertion, for each $n\geq3$ put
$H_n=(2\pi/n)I$ and $H_n'=H_n-2\pi P$, with $P$ a rank-one
diagonal projection. Formula \eqref{eq:gap}, with $T=2\pi$, gives
\[
 Q_4(H_n')-Q_4(H_n)=-4\pi^2(3+2/n)<0.
\]
Meanwhile $\|iH_n\|_{\op}=2\pi/n\to0$. For any proposed
radius $r>0$, choose $n>2\pi/r$ to obtain a counterexample with
operator norm less than $r$.
\end{proof}

In dimension one, write $F_c(i\theta)=a|\theta|$ with $a>0$.
A piecewise smooth circle curve from $1$ to $e^{i\theta}$ has a real
angle lift starting at $0$ and ending at $\theta+2\pi k$ for some integer
$k$. Its length is at least $a|\theta+2\pi k|\geq a|\theta|$ when
$|\theta|\leq\pi$; the curve $e^{it\theta}$ attains this bound.
Thus \eqref{eq:conjecture} holds in dimension one for this family.

On the trace-zero subspace, $F_c$ is a positive multiple of the
Hilbert--Schmidt norm. If $H$ is trace zero, the competitor
$e^{it(H-2\pi P)}$ has determinant $e^{-2\pi it}$ and does not stay in
$\mathrm{SU}(n)$. The winding construction therefore does not give a
competitor in the special-unitary setting of \cite[Theorem 4.1]{Larotonda}.
The results above do not assert an exact global distance for the
counterexamples, or a classification of all Ad-invariant norms.

\section{Coverage and verification}
\begin{center}\small
\begin{tabular}{@{}p{.48\linewidth}p{.45\linewidth}@{}}
\toprule
Source item & Coverage\\
\midrule
\cite{Larotonda}, Section 4.1, conjecture after Proposition 4.3
& Refuted (explicit counterexample), \eqref{eq:main-example} and Theorem~\ref{thm:family}.\\[3pt]
\cite{Larotonda}, Section 4.1, proposed input for spectral radius $\pi/2$
& Not asserted in the source; fails for the $\U(3)$ example in Corollary~\ref{cor:half}.\\[3pt]
Scalar threshold, fixed-dimension radius obstruction, and absence of a dimension-independent radius
& Not posed; obstructions proved in Corollaries~\ref{cor:scalar} and~\ref{cor:half}.\\[3pt]
\cite{Larotonda}, Theorem 2.1; trace family with $c=0$
& Prior proof cited: Antezana--Larotonda--Varela (2011/2014), \cite[Theorem 4.3]{ALV}; used in Theorem~\ref{thm:family}.\\[3pt]
Dimension-one trace family and trace-zero restriction
& Not posed; the circle case is proved above, and the trace term vanishes on $\mathfrak{su}(n)$.\\[3pt]
\cite{Larotonda}, Proposition 4.3; Theorems 1.1 and 4.1
& Not covered: segment convexity and the stated quotient theorems concern different assertions.\\[3pt]
\cite{Larotonda}, Remarks 5.4 and 5.5
& Not covered: uniqueness in quotients and infinite-dimensional extension.\\
\bottomrule
\end{tabular}
\end{center}

\paragraph{Verification scope.}
In the pinned
\href{https://github.com/the-omega-institute/trureturing/blob/a89e0fb9ba0f9002502ec8125455c2c2a9281443/D5/S3/QuadraticForms/LarotondaPrincipalLogDistanceRefutation.lean}{Lean source},
the declaration \texttt{result}, of type $\neg\,\texttt{claim}$, belongs to
the namespace
\begin{quote}\small\ttfamily
D5.S3.QuadraticForms.\\
LarotondaPrincipalLogDistanceRefutation.
\end{quote}
Here \texttt{claim} is the universal normalized Ad-invariant Finsler-norm
assertion \eqref{eq:conjecture} on skew-adjoint matrices, with piecewise
smooth unitary competitors. The declaration uses the
two-dimensional $c=4$ norm, endpoint and strict-inequality construction
in \eqref{eq:main-example}; it does not state the whole parameter-family
classification or Corollaries~\ref{cor:scalar} and~\ref{cor:half}.
Those conclusions have the written proofs above.
The accompanying \texttt{verify.py} gives exact rational finite checks
of the displayed coefficients and winding comparisons.

\begin{thebibliography}{9}
\bibitem{Larotonda}
G. Larotonda, \emph{Short curves in homogeneous spaces of the unitary
group for unitarily invariant norms}, arXiv:2610.31535v1 (2026),
\href{https://arxiv.org/abs/2610.31535v1}{arXiv:2610.31535v1}.

\bibitem{ALV}
J. Antezana, G. Larotonda and A. Varela,
\emph{Optimal paths for symmetric actions in the unitary group},
Communications in Mathematical Physics \textbf{328} (2014), 481--497,
\href{https://doi.org/10.1007/s00220-014-2041-x}{doi:10.1007/s00220-014-2041-x};
\href{https://arxiv.org/abs/1107.2439v1}{arXiv:1107.2439v1} (2011).

\bibitem{RS}
B. Russell and S. Stepney,
\emph{The geometry of speed limiting resources in physical models of computation},
International Journal of Foundations of Computer Science \textbf{28} (2017), 321--333,
\href{https://doi.org/10.1142/S0129054117500204}{doi:10.1142/S0129054117500204};
\href{https://arxiv.org/abs/1611.09220}{arXiv:1611.09220}.

\bibitem{LR}
G. Larotonda and I. Rey, \emph{Lie groups with a bi-invariant distance},
arXiv:2511.09699v2 (2025),
\href{https://arxiv.org/abs/2511.09699v2}{arXiv:2511.09699v2}.

\bibitem{LM}
G. Larotonda and M. Miglioli, \emph{Convexity of sums of eigenvalues of a
segment of unitaries}, arXiv:2408.16906v1 (2024),
\href{https://arxiv.org/abs/2408.16906v1}{arXiv:2408.16906v1}.

\bibitem{Metric}
G. Larotonda, \emph{Metric geometry of infinite-dimensional Lie groups
and their homogeneous spaces}, Forum Mathematicum \textbf{31} (2019),
1567--1605,
\href{https://doi.org/10.1515/forum-2019-0127}{doi:10.1515/forum-2019-0127};
\href{https://arxiv.org/abs/1805.02631v4}{arXiv:1805.02631v4}.
\end{thebibliography}
\paragraph{License.} \copyright\ 2026 Haobo Ma and Wenlin Zhang.
This work is licensed under
\href{https://creativecommons.org/licenses/by/4.0/}{Creative Commons Attribution 4.0 International (CC BY 4.0)}.
\end{document}
